\section{Fukaya 范畴（完整版本）}\label{sec:12}\label{section:12}
\subsection{Lagrangian 膜}\label{sec:12a}
设 $(M,\allowbreak \omega_M)$ 为满足 $2c_1(M)=0$ 的辛流形，配有相容殆复结构 $I_M$。选择二次复体积形式 $\eta_M^2$，即双典范丛 $K_M^2=\lambda_{\CC}^{\mathrm{top}}(TM)^{\otimes-2}$ 的一个处处非零截面；按假设，该丛平凡。令 $\alpha_M:\operatorname{Gr}(TM)\to S^1$ 为相应的平方相位映射，见\eref{11.34}。
\begin{definition*}
\emph{Lagrangian 膜}是三元组 $L^\#=(L,\allowbreak \alpha^\#,\allowbreak P^\#)$，由 Lagrangian 子流形 $L\subset M$ 及子丛 $TL\subset TM|L$ 上的膜结构 $(\alpha^\#,\allowbreak P^\#)$ 组成。具体而言，就是满足 $\exp(2\pi i\alpha^\#(x))=\alpha_M(TL_x)$ 的函数 $\alpha^\#:L\to\R$，即 $L$ 的分次，以及 $TL$ 的 Pin 结构 $P^\#$。
\end{definition*}
这种结构的基本理论是\secref{11j}一般讨论的特例。为与辛拓扑中更常用的概念联系，应指出：作为分次存在障碍的 Maslov 类 $\mu_L\in H^1(L)$，经边界映射 $H^1(L)\to H^2(M,\allowbreak L)$ 映为相对第一 Chern 类的两倍 $2c_1(M,\allowbreak L)$。另一方面，$\mu_L$ 的模二约化为 $w_1(L)-w_1(\xi_M|L)$，其中 $\xi_M\subset\lambda_{\CC}^{\mathrm{top}}(TM)$ 是使 $\pm\eta_M$ 取实值的实线丛。这与\stref{笔记}{11.18}密切相关：所有 Lagrangian 膜 $L^\#$ 都相对于 $\xi_M|L$ 典范定向。特别地，若 $\xi_M$ 平凡，等价地 $\eta_M^2$ 有全局平方根 $\eta_M$，则 Lagrangian 膜具有通常的定向；更确切地说，平方根的选择固定了这些定向。用 $L^\#\otimes\beta$ 表示以实线丛 $\beta\to L$ 改变 Pin 结构 $P^\#$ 的运算，以 $S^\sigma L^\#$ 表示逐点由\eref{11.36}给出的 $\sigma$ 重平移。\znote{174}{膜的定义把两件独立的几何数据拼在一起：$\alpha^\#$ 是平方相位映射的提升，它存在的充要条件是 Maslov 类 $\mu_L$ 为零；$P^\#$ 则是 $TL$ 上的 Pin 结构，用来处理 $TL$ 一般不可定向的情形。当 $\xi_M$ 平凡、从而每个 $L$ 可定向时，Pin 结构退化为通常的定向（至多差一个整体符号），这也是下文常直接用定向代替 Pin 数据的理由。分清这两项各管什么，才不会把“分次”与“定向”混为一谈。[来源：Seidel, Fukaya categories and Picard–Lefschetz theory, §11j 与例题 11.20。]}
\begin{note}\label{st:12.1}
Lagrangian 膜概念实质上只依赖 $\eta_M^2$ 的同伦类，有由 $H^1(M)$ 给出的实际选择自由度。此外，定义还有两个变体。首先，可固定类
\SourcePage{183}
$d\in H^1(M;\Z/2)$，用实线丛 $\delta$ 表示它，并将 $TL$ 上的 Pin 结构换成 $TL\oplus\delta|L$ 上的 Pin 结构。其次，给定 $w\in H^2(M;\Z/2)$，可考虑\stref{笔记}{11.16}意义下关于 $w|L$ 的扭曲 Pin 结构；也可将这两个想法结合。这里不继续展开，有兴趣的读者可参见\bcite{60}。不过，仍想概述如何将\secref{11h}的构造整体化，为这些可能性建立统一框架。

给定逐纤维分别限制为 Maslov 类与第二 Stiefel–Whitney 类的上同调类 $\mu_M\in H^1(\operatorname{Gr}(TM))$ 与 $w_M\in H^2(\operatorname{Gr}(TM);\Z/2)$，可如\eref{11.22}构造纤维化 $\operatorname{Gr}(TM)^\#\to\operatorname{Gr}(TM)$，其纤维为 $\Z\times\RP^\infty$。于是，抽象 Lagrangian 膜可定义为 Lagrangian 子流形 $L\subset M$，连同典范截面 $L\to\operatorname{Gr}(TM)|L$ 到 $\operatorname{Gr}(TM)^\#$ 的提升。为与此前术语联系，注意平方相位映射在二次体积形式 $\eta_M^2$ 的同伦类与上同调类 $\mu_M$ 之间建立双射。同样，最一般的类 $w_M$ 具有形式 $w_2(F)+w_1(F)\cdot\pi^*d+\pi^*w$，其中 $F\to\operatorname{Gr}(TM)$ 是万有丛，$\pi:\operatorname{Gr}(TM)\to M$ 是投影，$(d,\allowbreak w)$ 如前。\znote{175}{纤维 $\Z\times\RP^\infty$ 记录的正是膜结构的两部分：$\Z$ 分量给出 $\alpha^\#$ 的整值提升（$\alpha^\#$ 可整体加整数而不改变 $\alpha_M$），而 $\RP^\infty=K(\Z/2,1)$ 分类 Pin 结构，因为实线丛的同构类由 $H^1(\cdot;\Z/2)$ 刻画。所谓抽象膜，就是把 $L$ 的 Gauss 映射连同这两项额外拓扑数据一起提升为 $\operatorname{Gr}(TM)^\#$ 的截面；具体膜则对应于这些截面中满足相容条件的那些。注意 $\Z$ 分量并非多余：平移 $S^k$ 正是在这个因子上作加法，所以膜与它的平移在抽象膜空间里分居不同连通分支。[来源：Seidel, Fukaya categories and Picard–Lefschetz theory, §11h，尤见构造 (11.22)。]}
\end{note}

\subsection{模空间的定向}\label{sec:12b}
从现在起，回到通常的框架：$(M,\allowbreak \omega_M,\allowbreak \theta_M,\allowbreak I_M)$ 为带角正合辛流形，仍带有 $\eta_M^2$，且只考虑正合 Lagrangian 膜，即底层 Lagrangian 子流形正合的膜。

设 $(L_0^\#,\allowbreak L_1^\#)$ 为一对正合 Lagrangian 膜，$L_0,\allowbreak L_1$ 横截相交。在每个交点 $y\in L_0\cap L_1$ 的切空间 $(TM_y,\allowbreak \omega_{M,\allowbreak y},\allowbreak I_{M,\allowbreak y},\allowbreak \eta_{M,\allowbreak y}^2)$ 中，有显然的一对线性 Lagrangian 膜
\begin{equation}\label{eq:12.1}
\Lambda_{k,y}^\#=(\Lambda_{k,y}=(TL_k)_y,\ \alpha_{k,y}^\#=\alpha_k^\#(y),\ P_{k,y}^\#=(P_k^\#)_y).
\end{equation}
取\eref{11.25}定义的指标与定向空间，记为 $i(y)=i(\Lambda_{0,\allowbreak y}^\#,\allowbreak \Lambda_{1,\allowbreak y}^\#)$、$o(y)=o(\Lambda_{0,\allowbreak y}^\#,\allowbreak \Lambda_{1,\allowbreak y}^\#)$。更一般地，允许任意 $(L_0,\allowbreak L_1)$，选择 Floer 数据 $(H,\allowbreak J)$，并考虑 $y\in\mathcal C(L_0,\allowbreak L_1)$。此时\eref{12.1}中的 $\Lambda_{k,\allowbreak y}^\#$ 起初位于不同向量空间 $TM_{y(k)}$ 中；不过，与\secref{11l}的讨论平行，可利用 $H$ 的 Hamilton “流”在 $y$ 附近的线性化，结合 $\operatorname{Gr}(TM_{y(t)})$ 上的平方相位函数，将 $\Lambda_{0,\allowbreak y}^\#$ 移到 $TM_{y(1)}$ 中。然后如前定义 $i(y)$ 与 $o(y)$。由\eref{11.39}，也可将它们视为任意选取的、对应于 $(\Lambda_{0,\allowbreak y}^\#,\allowbreak \Lambda_{1,\allowbreak y}^\#)$ 的定向算子的指标与行列式线，将这种算子记为 $D_{H,\allowbreak y}$。

假设 $S$ 为配有条带状端及膜标记 $L_0^\#,\allowbreak \ldots,\allowbreak L_d^\#$ 的 $(d+1)$ 点圆盘；膜标记就是连同膜结构的 Lagrangian 标记。为 $(L_0,\allowbreak L_d)$ 及 $(L_{k-1},\allowbreak L_k)$（$k=1,\allowbreak \ldots,\allowbreak d$）选择 Floer 数据，并在 $S$ 上选择相容
\SourcePage{184}
扰动数据。在正则点 $u\in\M_S(y_0,\allowbreak \ldots,\allowbreak y_d)$ 处，有
\begin{equation}\label{eq:12.2}
\begin{aligned}
\dim_u\M_S(y_0,\ldots,y_d)&=i(y_0)-i(y_1)-\cdots-i(y_d),\\
(\lambda^{\mathrm{top}}T\M_S(y_0,\ldots,y_d))_u
&\cong\lambda^{\mathrm{top}}\ker(D_{S,u})\\
&\cong o(y_0)\otimes o(y_1)^\vee\otimes\cdots\otimes o(y_d)^\vee.
\end{aligned}
\end{equation}
第二式的意思是：模空间相对于各端的定向空间典范定向。\eref{12.2}是\stref{命题}{11.13}的推论，不过，直接重述相关粘合论证也许更有启发。为 $k=1,\allowbreak \ldots,\allowbreak d$ 选择定向算子 $D_{H,\allowbreak y_k}$。依定义，它们是 $H$ 上的 $\bar\partial$ 算子，大致说来，其边界条件子丛带有膜结构。线性化算子 $D_{S,\allowbreak u}$ 带类似数据。因此，依次将 $D_{S,\allowbreak u}$ 与 $D_{H,\allowbreak y_d},\allowbreak \ldots,\allowbreak D_{H,\allowbreak y_1}$ 粘合，并与给定膜结构相容，便得到 $(\Lambda_{0,\allowbreak y_0}^\#,\allowbreak \Lambda_{1,\allowbreak y_0}^\#)$ 的一个定向算子，记为 $D_{H,\allowbreak y_0}$。粘合公式与\eref{11.39}立即给出
\begin{equation}\label{eq:12.3}
\begin{aligned}
i(y_0)&=\operatorname{index}(D_{S,u})+i(y_1)+\cdots+i(y_d),\\
o(y_0)&\cong\det(D_{S,u})\otimes o(y_d)\otimes\cdots\otimes o(y_1).
\end{aligned}
\end{equation}
正则情形下，$D_{S,\allowbreak u}$ 满射，且 $\ker(D_{S,\allowbreak u})\cong T\M_S(y_0,\allowbreak \ldots,\allowbreak y_d)_u$，故得\eref{12.2}。
\begin{note}\label{st:12.2}
更精确地，这里的约定是对偶带有典范配对 $o(y_k)^\vee\otimes o(y_k)\to\R$。特别地，从\eref{12.2}中的同构出发，在右侧按递减次序与所有 $o(y_k)$ 作张量积，就还原出\eref{12.3}。
\end{note}

需要知道，这些定向在粘合非齐次伪全纯映射这一基本运算下如何变化。取膜标记为 $(L_0^\#,\allowbreak \ldots,\allowbreak L_n^\#,\allowbreak L_{n+m}^\#,\allowbreak \ldots,\allowbreak L_d^\#)$ 的 $(d+2-m)$ 点圆盘 $S_1$，以及标记为 $(L_n^\#,\allowbreak \ldots,\allowbreak L_{n+m}^\#)$ 的 $(m+1)$ 点圆盘 $S_2$。以某个粘合长度 $l$，将 $S_1$ 的第 $n+1$ 个出向端与 $S_2$ 的入向端粘合，得到继承标记 $(L_0^\#,\allowbreak \ldots,\allowbreak L_d^\#)$ 的 $(d+1)$ 点圆盘 $S$。为所有相关的 Lagrangian 子流形对选择 Floer 数据，在 $S_1,\allowbreak S_2$ 上选择相容扰动数据，并诱导 $S$ 上的扰动数据。假设对某个 $\tilde y$，$u_1\in\M_{S_1}(y_0,\allowbreak \ldots,\allowbreak y_n,\allowbreak \tilde y,\allowbreak y_{n+m+1},\allowbreak \ldots,\allowbreak y_d)$ 与 $u_2\in\M_{S_2}(\tilde y,\allowbreak y_{n+1},\allowbreak \ldots,\allowbreak y_{n+m})$ 是各自模空间的正则点。只要 $l$ 足够大，先用单位分解直接拼接 $u_k$，再应用隐函数定理，粘合过程就产生正则点 $u\in\M_S(y_0,\allowbreak \ldots,\allowbreak y_d)$。这一过程延拓到模空间的邻近点，实际上给出 $(u_1,\allowbreak u_2)$ 与 $u$ 的邻域之间的局部微分同胚 $g$。即使从这段粗略描述中，也应能理解其微分 $Dg$ 近似于线性化算子的\eref{11.9}
\SourcePage{185}
所给映射。为得出所需结论，考虑下图：
\SourceEquationDiagram{diagram-12-4}{0.65423006}{12.4}
水平同构是\eref{12.2}，右侧向上箭头是 $o(\tilde y)$ 与其对偶之间的典范配对。则\eref{12.4}在相差 Koszul 约定规定的符号的意义下交换，该符号为 $(-1)^\Delta$，其中
\[
\Delta=\left(\sum_{k>n+m}i(y_k)\right)\operatorname{index}(D_{S_2,u_2}).
\]
此符号的来源可追溯到行列式线的定义，见\secref{11a}，尤其是\eref{11.2}。\znote{177}{这个 Koszul 符号是记账方式而非笔误：行列式线是分次线，两个奇次因子交换次序会带来 $-1$。粘合公式 (11.9) 把 $D_{S,u}$ 与各 $D_{H,y_k}$ 依次拼接，每交换一次记一个符号，累积起来即 $(-1)^\Delta$，其中 $\Delta$ 由外侧交点的指标之和与 $D_{S_2,u_2}$ 的指标相乘得到。把这一条弄懂，12g 中 $\ddagger$、$\diamond$、$\dagger$ 三个符号就有了统一的来源，不必逐条死记。读者不妨按 (12.22) 自己推一遍 $\ddagger=m(d-n)+m+n$，看每一项来自哪一次交换。[来源：Seidel, Fukaya categories and Picard–Lefschetz theory, §11a–11b，尤见 (11.2) 与命题 11.13。]}

上述讨论可以沿两个方向推广。为叙述方便，分别处理这两个方向，将二者结合这一直接的工作留给读者。

\subsection{更一般的 Riemann 曲面}\label{sec:12c}
第一个想法是扩大 Riemann 曲面 $S$ 的范围。先看紧化 $\hat S$ 的每个边界圆周都恰有一个入向无穷远点的情形。对适当模空间 $\M_S(\{y_\zeta\})$ 的正则点 $u$，\stref{命题}{11.13}表明
\[
\begin{aligned}
\dim_u\M_S(\{y_\zeta\})&=n(\chi(S)-|\Sigma^-|)+\sum_{\zeta\in\Sigma^\pm}\mp i(y_\zeta),\\
(\lambda^{\mathrm{top}}T\M_S(\{y_\zeta\}))_u&\cong\bigotimes_{\zeta\in\Sigma^\pm}o(y_\zeta)^{\mp1},
\end{aligned}
\]
而\eref{12.4}的推广没有困难。只要每个边界圆周包含奇数个入向无穷远点，同样成立；不过，粘合时追踪符号行为更棘手，因为\stref{命题}{11.13}的相关部分多了一个步骤。

余下情形稍不直接。对 $\hat S$ 的每个包含偶数个入向无穷远点的边界圆周，选择属于该圆周的 $\partial S$ 的一个分支，再在该分支上选一个标记点。将这些选择记为 $(C_j,\allowbreak z_j)$，再运用\stref{命题}{11.13}得到
\begin{equation}\label{eq:12.5}
(\lambda^{\mathrm{top}}T\M_S(\{y_\zeta\}))_u
\cong\bigotimes_j(\lambda^{\mathrm{top}}TL_{C_j})_{u(z_j)}\otimes\bigotimes_{\zeta\in\Sigma^\pm}o(y_\zeta)^{\mp1}.
\end{equation}
消除额外因子的显然办法是确保 $L_{C_j}$ 有定向，例如假定 $\xi_M$ 平凡。另一种技巧是将额外因子吸收到定向空间的改变中。为使记号简单，
\SourcePage{186}
只说明一个特殊情形：$S$ 为具有 $e$ 个入向和 $d$ 个出向无穷远点的圆盘，其中 $e$ 为偶数，$d>0$，排列如\sr{fig:12.1}{图~12.1}。将涉及的 Lagrangian 膜记为 $L_k^{\pm,\allowbreak \#}$，伪全纯映射 $u$ 的极限记为 $y_k^\pm$。将标记点 $z$ 向 $\zeta_d^+$ 移动，使 $(TL_d^+)_{u(z)}$ 渐近等于 $\Lambda_{y_d^+,\allowbreak 1}$，因而可识别
\begin{equation}\label{eq:12.6}
\begin{aligned}
&(\lambda^{\mathrm{top}}TL_d^+)_{u(z)}\otimes o(\Lambda_{y_d^+,0}^\#,\Lambda_{y_d^+,1}^\#)\\
&\qquad\cong o(\Lambda_{y_d^+,0}^\#,\Lambda_{y_d^+,1}^\#\otimes\lambda^{\mathrm{top}}\Lambda_{y_d^+,1}).
\end{aligned}
\end{equation}
据此，\eref{12.5}变为
\begin{equation}\label{eq:12.7}
\begin{aligned}
&\lambda^{\mathrm{top}}T\M_S(y_1^-,\ldots,y_e^-,y_1^+,\ldots,y_d^+)\\
&\quad\cong o(\Lambda_{y_e^-,0}^\#,\Lambda_{y_e^-,1}^\#)\otimes\cdots\otimes o(\Lambda_{y_1^-,0}^\#,\Lambda_{y_1^-,1}^\#)\\
&\qquad\otimes o(\Lambda_{y_1^+,0}^\#,\Lambda_{y_1^+,1}^\#)^\vee\otimes\cdots
\otimes o(\Lambda_{y_d^+,0}^\#,\Lambda_{y_d^+,1}^\#\otimes\lambda^{\mathrm{top}}\Lambda_{y_d^+,1})^\vee.
\end{aligned}
\end{equation}
\SourceFigure{figure-12-1}{0.30323044}{12.1}

\subsection{带标记点圆盘的族}\label{sec:12d}
沿另一个方向，取膜标记为 $(L_0^\#,\allowbreak \ldots,\allowbreak L_d^\#)$ 的 $(d+1)$ 点圆盘族 $\mathcal S\to\RG$。选择扰动数据，考虑相应的一个参数化模空间 $\M_{\mathcal S}(y_0,\allowbreak \ldots,\allowbreak y_d)$。正则点 $(r,\allowbreak u)$ 处的切空间现在由\eref{9.18}中的扩展算子 $D_{\mathcal S,\allowbreak r,\allowbreak u}$ 的核给出。将该算子的第一个有限维分量变形成零，得到
\begin{equation}\label{eq:12.8}
\begin{aligned}
\dim_{(r,u)}\M_{\mathcal S}(y_0,\ldots,y_d)&=\dim_r\RG+i(y_0)-i(y_1)-\cdots-i(y_d),\\
(\lambda^{\mathrm{top}}T\M_{\mathcal S}(y_0,\ldots,y_d))_{(r,u)}
&\cong(\lambda^{\mathrm{top}}T\RG)_r\otimes o(y_0)\otimes o(y_1)^\vee\otimes\cdots\otimes o(y_d)^\vee.
\end{aligned}
\end{equation}
假设有两个族 $\mathcal S_k\to\RG_k$（$k=1,\allowbreak 2$），如前分别标为 $(L_0^\#,\allowbreak \ldots,\allowbreak L_n^\#,\allowbreak L_{n+m}^\#,\allowbreak \ldots,\allowbreak L_d^\#)$ 与 $(L_n^\#,\allowbreak \ldots,\allowbreak L_{n+m}^\#)$，并配有适当扰动数据。小节 \sr{sec:9e}{(9e)} 与
\SourcePage{187}
\sr{sec:9i}{(9i)} 所述的族粘合过程，在两顶点树这一特殊情形下，给出诱导族 $\mathcal S\to\RG$，其中 $\RG=(-1;0)\times\RG_1\times\RG_2$。对应的非齐次伪全纯映射粘合，取正则点 $(r_1,\allowbreak u_1)\in\M_{\mathcal S_1}(y_0,\allowbreak \ldots,\allowbreak y_n,\allowbreak \tilde y,\allowbreak y_{n+m+1},\allowbreak \ldots,\allowbreak y_d)$ 与 $(r_2,\allowbreak u_2)\in\M_{\mathcal S_2}(\tilde y,\allowbreak y_{n+1},\allowbreak \ldots,\allowbreak y_{n+m})$，以及足够接近零的粘合参数 $\rho\in(-1;0)$，产生元素 $(r,\allowbreak u)\in\M_{\mathcal S}(y_0,\allowbreak \ldots,\allowbreak y_d)$。此外，它延拓为 $(-1;0)\times\M_{\mathcal S_1}\times\M_{\mathcal S_2}$ 与 $\M_{\mathcal S}$ 之间的局部微分同胚 $g$。注意，由于其中用到了隐函数定理，$r$ 通常不等于 $(\rho,\allowbreak r_1,\allowbreak r_2)$。不过两点很接近；更精确地，$\rho\to0$ 时两者渐近接近。特别地，可识别
\begin{equation}\label{eq:12.9}
(\lambda^{\mathrm{top}}T\RG)_r\cong\R\otimes(\lambda^{\mathrm{top}}T\RG_1)_{r_1}\otimes(\lambda^{\mathrm{top}}T\RG_2)_{r_2}.
\end{equation}
\eref{12.4}的适当推广是下图：
\SourceEquationDiagram{diagram-12-10}{0.64176854}{12.10}
右侧向上箭头结合了配对 $o(\tilde y)^\vee\otimes o(\tilde y)\cong\R$ 与\eref{12.9}。该图在相差 $(-1)^\diamond$ 的意义下交换，其中
\[
\diamond=\dim(\RG_2)\operatorname{index}(D_{S_1,u_1})+
\left(\sum_{k>n+m}i(y_k)\right)\operatorname{index}(D_{S_2,u_2}).
\]
这里的指标是曲面 $S_k=(\mathcal S_k)_{r_k}$ 上通常线性化算子的指标，并非涉及参数空间的扩展算子的指标。
